\vspace*{0.6\baselineskip}
\begin{center}
  {\Huge\bfseries Introduction générale}
\end{center}
\addcontentsline{toc}{chapter}{Introduction générale}
\markboth{Introduction générale}{}
\vspace{1.5cm}

\textbf{Poulina Group Holding (PGH)} gère une centaine de filiales industrielles en Tunisie et génère chaque jour des volumes de livraisons que ses équipes planifient encore manuellement : un seul planificateur, des affectations par téléphone, des itinéraires non tracés et des paiements sans reçu numérique. Ce mode de fonctionnement tient jusqu'à une certaine échelle ; au-delà, les erreurs s'accumulent et les coûts logistiques deviennent difficiles à maîtriser.\\

Notre projet de fin d'études, réalisé au sein de \textbf{GREPSYS Tunisia}, répond à ce besoin concret. Nous avons conçu et développé \textbf{Fleet~AI}, une application de gestion de flotte intégrée à la plateforme \textbf{WFMS} (Workflow Management System) déjà en production chez PGH. L'application couvre le cycle logistique complet : référentiels de ressources, gestion des commandes, optimisation des itinéraires, suivi GPS, gestion des paiements, prédiction des retards et tableaux de bord décisionnels. Elle n'est pas un remplacement du WFMS existant --- elle l'augmente.\\

Ce rapport suit le déroulement réel du projet :

\begin{itemize}
    \item[$\bullet$] Le \textbf{premier chapitre} présente GREPSYS, PGH et le cahier des charges complet de Fleet~AI.
    \item[$\bullet$] Le \textbf{deuxième chapitre} analyse les solutions de gestion de flotte existantes et explique pourquoi aucune ne convient au contexte de PGH.
    \item[$\bullet$] Le \textbf{troisième chapitre} décrit la méthode Scrum retenue, l'environnement technique et l'architecture logicielle.
    \item[$\bullet$] Les \textbf{chapitres 4 à 8} détaillent chacun un sprint : backlog, conception, réalisation et tests.
\end{itemize}
